\chapter{Background on Nonlinear Optimization}
In this chapter we will go through the necessary background to understand
optimality conditions for nonlinear constrained optimization problems. We
will introduce necessary optimality conditions. In the case we have a 'nice
enough', convex set we can show that under certain qualifications the
neccessary conditions are also sufficient for optimality. These are all very
well known results from the theory of nonlinear optimization, hence this
chapter is based on the lecture notes in
\cite{bot_optimization}.
\section{First Order Optimality Conditions}
Generally we are considering a nonempty subset $X \subseteq \mathbb{R}^{n}$
for $n \in \mathbb{N}$ and a function $f: \mathbb{R}^{n}  \to \mathbb{R}$ in
the context of an optimization problem
\begin{align}
    \min_{x \in X}\; &f(x) \label{eq: basic_opt}.
\end{align}
A local minimum of this problem is an element $x^{*} \in X$,
such that there is an $\varepsilon > 0$ ball, wrt. to the euclidean norm,
i.e.
\begin{align}
    U_\varepsilon(x^{*}) = \{x \in X: \|x - x^{*}\| < \varepsilon\},
\end{align}
such that
\begin{align}
    f(x^{*} ) \le f(x) \qquad \forall x \in U_\varepsilon(x^{*})\cap X.
\end{align}
A global minimum is an element $x^{*} \in X$ such that the inequality holds
on the whole set $X$.
\begin{align}
    f(x^{*}) \le f(x) \qquad \forall x \in X.
\end{align}
To better understand local minima we look at directions, from a given point
$x_0$, from where we can move and still remain in $X$. In this sense we
introduce a cone
\begin{align}
    T_X(x_0) := \left\{ d \in \mathbb{R}^{n}\Big|\;\exists (x^k)_{k \in
    \mathbb{N}} \subseteq X,\; \exists (t_k)_{k \in \mathbb{N} }\searrow
0\;\; : \lim_{k \to \infty} \frac{x^k-x_0}{t_k} = d  \right\},
\end{align}
called the Bouligard tangent cone. It is nonempty since $0 \in T_X(x_0)$ for
the choice of $x^{k} = x_0$ and $t_k = \frac{1}{k}$. It is really a cone
since it is scaling invariant for $d \in T_X(x_0)$ also $\lambda d \in
T_X(x_0)$, for all $\lambda > 0$, by multiplying $(t_k)_{k \in \mathbb{N} }$
with $\frac{1}{\lambda}$.
\begin{proposition}
    \label{prop: descent}
    let $x^{*}$ be a local minimum of \ref{eq: basic_opt}, where $f \in
    C^1(U_{\varepsilon}(x^{*}))$, for an $\varepsilon >0$. Then
    \begin{align}
        \nabla f(x^{*})^{T}d \ge 0 \qquad \forall d \in T_X(x^{*}).
    \end{align}
\end{proposition}
\begin{proof}
Fix $d \in T_X(x^{*})$, then there are sequences $(x^{k})_{k \in \mathbb{N} }$
and $(t_k)_{k \in \mathbb{N}}$ such that
\begin{align}
    \lim_{k \to \infty} \frac{x^{k}-x^{*}}{t_k} = d.
\end{align}
Multiplying by $t_k$ gives
\begin{align}
   \lim_{k \to \infty} x^{k} - x^{*} =
   \lim_{k \to \infty}  \frac{x^{k}-x^{*}  }{t_k}t_k =
   \lim_{k \to \infty} dt_k = 0.
\end{align}
Since $x^{*}$ is a local minimum of \ref{eq: basic_opt} we have that
$f(x^{*}) \le f(x)$ for all $x \in U_\varepsilon(x^{*}) \cap X$. So there is
a $k_\varepsilon \in \mathbb{N}$ such that $x^{k} \in U_\varepsilon(x^{*})$
for all $k\ge k_\varepsilon$. By the mean value theorem there is a $\xi^{k}
\in \{\lambda x^{*} + (1-\lambda)x^{k}| \lambda \in [0, 1]  \} $, such that
\begin{align}
    f(x^{k}) - f(x^{*} ) = \nabla f(\xi^{k} )(x^{k} -x^{*} ).
\end{align}
Let $\xi^{k} = \lambda_k x^{*} + (1-\lambda_k)x^{k}$, then for all $k \ge
k_\varepsilon$ we have
\begin{align}
    \|\xi^{k} -x^{k} \| &= \|\lambda_k x^{*} + (1-\lambda_k) x^{k}
    -x^{*}\| \\
                &=|\lambda_k-1| \| x^{k}  - x^{*} \|\\
                &\le \| x^{k}  - x^{*} \| \to 0.
\end{align}
And thus $\xi^{k} \to x^{*}$. By the continuity of the gradient in the
neighborhood of $x^{*}$ we also have $\nabla f(\xi^{k}) \to \nabla f(x^{*})$,
thereby
\begin{align}
    \nabla f(x^{*})^{T} d &= \lim_{k \to \infty} \frac{\nabla f(\xi^{*})
    (x^{k}-x^{*})}{t_k} \\
              &=\lim_{k \to \infty} \frac{f(x^{k}) - f(x^{*})  }{t_{k} } \\
              &\ge 0.
\end{align}
\end{proof}
The statement in Proposition \ref{prop: descent}, basically implies that at a
local minimum $x^{*}$ there is no tangent direction or `descent direction'
for which $\phi '(0) < 0$, where $\phi(t) := f(x^{*} + t d)$.

Bringing everything together gives the first order optimality condition for
unconstrained problems.
\begin{theorem}
    Let $x^{*}$ be a local minimum of \ref{eq: basic_opt} and $f \in C^{1}
    (U_\varepsilon(x^{*} ))$ for an $\varepsilon >0$. Then $x^{*}$ is a
    critical point of $f$, i.e.
    \begin{align}
        \nabla f(x^{*}) = 0.
    \end{align}
\end{theorem}
\begin{proof}
    Since $x^{*}$ is a local minimum of \ref{eq: basic_opt} and $f \in C^{1}
    (U_\varepsilon(x^{*} ))$ then $x^{*} $ is also a local minimum in of f in
    $U_\varepsilon(x^{*})$. By Proposition \ref{prop: descent} we have that
    $\nabla f(x^{*} )^{T}  d \ge 0$, for all $d \in T_{U_\varepsilon(x^{*})
    }(x^{*} ) = \mathbb{R} ^{n}$. Fix $d \neq 0$ and plug in $d$ and $-d$,
    giving $\nabla f(x^{*})^{T} d = 0$. Hence $\nabla f(x^{*} ) = 0$.
\end{proof}
From now on consider a particular set $X$, the feasible set, defining a
constrained optimization problem through
\begin{align}
    X = \left\{ x \in \mathbb{R}^{n}\Big|
        \begin{array}{l}
            g_i(x) \le 0,\; i=1,\ldots,m\\
            h_j(x) = 0,\; j=1,\ldots,p
    \end{array} \right\}.
\end{align}
Then we have a general constrained nonlinear optimization problem
\begin{align}
    \min\;&f(x) \label{eq: gen_constraned}\\
    \text{s.t.:}\;&g_i(x) \le 0,\; i=1,\ldots,m \nonumber\\
        &h_j(x) = 0,\; j=1,\ldots,p \nonumber\\
        &x \in \mathbb{R}^{n}\nonumber,
\end{align}
i.e. $x \in X$, in the feasible set.

Generally the Bouligard tangent cone $T_X(x_0)$ for a $x_0 \in X$ is not a
practical object, since it cannot be determined out of the box. Hence, we
introduce the linearized tangent cone, which will serve as a replacement and
can be easily determined for a given feasible set $X$. Also, in a specific
setting the linearized tangent cone will be valid representation of the
Bouligard tangent cone. For an element $x_0 \in X$ consider
\begin{align}
    \mathcal{A}(x_0) = \left\{ i = 1,\ldots,m|\;g(x_0) = 0 \right\},
\end{align}
the set of active indices at $x_0$ and
\begin{align}
    \mathcal{I}(x_0) = \left\{1,\ldots,m \right\}\setminus\mathcal{A}(x_0),
\end{align}
the set of inactive indices at $x_0$. The linearized tangent cone of $X$ at
$x_0$ is defined as
\begin{align}
    T_{\text{lin} }(x_0) = \left\{
        d\in \mathbb{R}^{n} \Big|
        \begin{array}{l}
            \nabla g_i(x_0)^{T} d \le 0,\; i \in \mathcal{A}(x_0)\\
            \nabla h_j(x_0)^{T} d = 0,\; j=1,\ldots,p.
        \end{array} \right\}.
\end{align}
The linearized tangent cone is really a cone since the scaling with
$\lambda>0$ does not change the outcome of the equalities and inequalities.
\begin{lemma}
    Let $x_{0} \in X$, then it holds that $T_X(x_0) \subseteq
    T_{\text{lin}}(x_0)$.
\end{lemma}
\begin{proof}
    We start with a $x_0 \in X$ and $d \in T_X(x_0)$. Then there are
    sequences by the definition of the Bouligard tangent cone, $(x^{k})_{k\in N}$ and
    $(t_k)_{k \in N} \searrow 0 $, such that
    \begin{align}
        \lim_{ k \to \infty} \frac{x^k-x_0}{t_k} = d.
    \end{align}
    First we prove that $\nabla g_i(x_0)^{T} d \le 0$, for all active indices
    $i \in \mathcal{A}(x_0)$. For this we use the mean value theorem, for
    all $k\ge1$ there is a $\xi^{k}  \in \{\lambda x_0 + (1-\lambda) x^{k}|\;
    \lambda \in [0, 1]\} $ such that
    \begin{align}
        g_i(x^{k}) - g_i(x_0) = \nabla g_i(\xi^{k})^{T} (x^{k} - x_0).
    \end{align}
    Since $i$ is active, we have $g_i(x_0) = 0$ and thus
    \begin{align}
         \nabla g_i(\xi^{k})^{T} (x^{k} - x_0) = g_i(x^{k}) \le 0.
    \end{align}
    Dividing by $t_k$ and passing the limit gives
    \begin{align}
        0 \ge \lim_{k \to \infty} \nabla g_i(\xi^{k})^{T}  \frac{x^k-x_0}{t_k}
        = \nabla g_i(x_0)^{T} d.
    \end{align}
    For the second statement in $T_{\text{lin}}(x_0)$ we need to show that
    $\nabla h_j(x_0)^{T} d = 0$ for all $j=1,\ldots,p$. For this we can use
    the same argumentation, by the mean value theorem for
    all $k\ge1$ there is a $\xi^{k}  \in \{\lambda x_0 + (1-\lambda) x^{k}|\;
    \lambda \in [0, 1]\} $ such that
    \begin{align}
        h_j(x^{k}) - h_j(x_0) = \nabla h_j(\xi^{k})^{T} (x^{k} - x_0).
    \end{align}
    Since $h_j(x_0) = 0$ and $h_j(x^{k}) = 0$  we have that
    \begin{align}
        \nabla h_j(\xi^{k})^{T} (x^{k} - x_0) = 0.
    \end{align}
    Dividing by $t_k$ and passing the limit gives
    \begin{align}
        0 = \lim_{ k \to \infty} \nabla h_j(\xi^{k})^{T} \frac{x^{k} - x_0}
        {t_k}= \nabla h_j(x_0)^{T} d.
    \end{align}
    Thus $d \in T_{\text{lin}}(x_0)$.
\end{proof}

\begin{lemma}[Lemma of Ferkas]
    \label{lem: ferkas}
    Let $A \in \mathbb{R}^{m \times n}$ and $b \in \mathbb{R}^{m}$. Then
    the two statements are equivalent
    \begin{enumerate}[nosep]
        \item The system $Ax = b$ has a solution $x\ge0$.
        \item For all $d \in \mathbb{R}^{m}$ with $A^{T} d \ge 0$ it holds
            that $b^{T} d  \ge 0$.
    \end{enumerate}
\end{lemma}
%\begin{proof}
%    We can use the strong duality result to prove equivalence.
%    For (1) $\Rightarrow$ (2), consider
%    \begin{align}
%        \text{min}\; & 0^{T} x\\
%        \text{s.t.:}\;& Ax = b \nonumber \\
%            & x \ge 0. \nonumber
%    \end{align}
%    Thus the dual of wrt. $d \in \mathbb{R}^{m}$ is
%    \begin{align}
%        \text{max}\; & b^{T}d \\
%        \text{s.t.:}\;& Ad \ge 0 \nonumber
%    \end{align}
%    By strong duality, the optimal value of the dual is $0 \in
%    \mathbb{R}^{m}$. Hence, for all $d$ with $A^{T}d \ge 0$ we have that
%    $b^{T}d \ge 0$.
%
%    For (2) $\Rightarrow$ (1), we have that $b^{T} d \ge 0$ for all $d$ with
%    $A^{T} d \ge 0$. By strong duality the optimal value of the primal is
%    $0$, and by the dual the optimal solution needs to be $x\ge 0$. By
%    feasibility of the primal there is an $x\ge0$ satisfying $Ax=b$.
%\end{proof}
\begin{theorem}
    \label{thm: duality}
    Let $x_0 \in X$, the it holds that
    \begin{align}
        -(T_{lin}(x_0))^{*} = N_{\text{lin} }(x_0).
    \end{align}
    where
    \begin{align}
        N_{\text{lin}} =
        \left\{
            \sum_{i=1}^{m} \lambda_i \nabla g_i(x_0)
            + \sum_{j=1}^{p} \mu_j\nabla h_j(x_0)
             \left|
             \begin{array}{l}
                \lambda_i \ge 0,\; i \in \mathcal{A}(x_0),\\
                \lambda_i = 0,\; i \in \mathcal{I}(x_0),\\
                \mu_j \in \mathbb{R},\; j=1,\ldots,p
            \end{array}
        \right\} \right.
    \end{align}
\end{theorem}
\begin{proof}
    Direction $\supseteq$: Let $s \in N_{\text{lin} }(x_0)$ and $d \in
    T_{\text{lin} }(x_0)$. Then
    \begin{align}
        -s^{T} d &= -\sum_{i=1}^{m} \lambda_i\nabla g_i(x_0)^{T}d
        -\sum_{j=1}^{p} \mu_j \nabla h_j(x_0)^{T} d\\
                 &\ge 0.
    \end{align}
    by $\lambda_i \ge0$, $\nabla g_i(x_0)^T d \le 0$ and $\nabla
    h_j(x_0)^{T} d =0 $. Hence $s \in -(T_{\text{lin} }(x_0))^{*}$.

    Direction $\subseteq$: Let $s \in -(T_{\text{lin} }(x_0))^{*}$, then by
    definition $-s^{T} d \ge 0$ for all $d \in T_{\text{lin}}(x_0)$. We can
    use Lemma \ref{lem: ferkas} with the matrix $A \in \mathbb{R}^{n\times
    (r + 2p)}$, with active indices $\mathcal{A}(x_0) = \{i_1, \ldots, i_r\}
    $ defined as
    \begin{align}
        A:=\begin{pmatrix} A_g, A_h, -A_h
        \end{pmatrix},
    \end{align}
    where
    \begin{align}
        A_g &=
        \begin{pmatrix}
        -\nabla g_{i_1}(x_0),\ldots,-\nabla g_{i_r}(x_0)
        \end{pmatrix},\\
        A_h &=
        \begin{pmatrix}
            \nabla h_1(x_0),\ldots,\nabla h_p(x_0)
        \end{pmatrix} .
    \end{align}
    By Lemma \ref{lem: ferkas} there is a $\nu = (\lambda, \mu_1, \mu_2) \in
    \mathbb{R}_+^{r}\times \mathbb{R}^{p}_+\times \mathbb{R}_+^{p}$
    satisfying
    \begin{align}
        A\nu = -s.
    \end{align}
    This is equivalent to
    \begin{align}
        \sum_{i=1}^{m} -\lambda_i\nabla g_i(x_0)
        +\sum_{j=1}^{p} (\mu_1)_j\nabla h_j(x_0)
        +\sum_{j=1}^{p} -(\mu_2)_j\nabla h_j(x_0) = -s.
    \end{align}
    By setting $\mu = \mu_2 - \mu_1$ and the fact that $\lambda_i \ge 0$ for
    all $i \in \mathcal{A}(x_0)$ and $ \lambda_i =0$ for all $i \in
    \mathcal{I}(x_0)$ we have
    \begin{align}
       s =  \sum_{i=1}^{m} \lambda_i\nabla g_i(x_0)
        +\sum_{j=1}^{p} \mu_j\nabla h_j(x_0),
    \end{align}
    thus $s \in N_{\text{lin}}(x_0)$.
\end{proof}

Next we will introduce the KKT conditions
\begin{definition}
    \label{def: KKT}
    The function $L: \mathbb{R}^{n} \times \mathbb{R}^{m} \times
    \mathbb{R}^{p} \to \mathbb{R}$ defined as
    \begin{align}
        L(x, \lambda, \mu)= f(x) + \lambda^{T} g(x) + \mu^{T} h(x)
    \end{align}
    is called the Lagrangian function of the optimization problem
    \ref{eq: gen_constraned}. The conditions
    \begin{align}
        \begin{cases}
            \nabla_x L(x, \lambda, \mu) = 0,\\
            \lambda \ge 0,\; g(x) \le 0,\; \lambda^{T} g(x) = 0\\
            h(x) = 0
        \end{cases}
    \end{align}
    are called KKT optimality conditions, where
    \begin{align}
        \nabla_x L(x, \lambda, \mu) &= \nabla f(x) + \lambda^{T} \nabla
            g(x) +\mu^{T} \nabla h(x)\\
        &= \nabla f(x) + \sum_{i=1}^{m} \lambda_i \nabla g_i(x)
            + \sum_{j=1}^{p} \mu_j \nabla h_j(x).
    \end{align}
    An element $(x^{*}, \lambda^{*}, \mu^{*})$ fulfilling  the KKT optimality conditions is called a KKT-point. The vectors $\lambda^{*}$ and $\mu^{*}$ are
    called Lagrange multipliers associated with the corresponding restrictions,
    i.e. $g(x^{*}) \le 0$ and $h(x^{*}) = 0$.
\end{definition}
In the case there are no restrictions, the KKT optimality conditions are
$\nabla f(x) = 0$. Thus the scope of these conditions is a generalization of
the first order optimality for constrained optimization and these play are
necessary conditions for a solution to be optimal. The KKT conditions
together with Abadie constraint qualification are necessary for a optimality.
However, in the case the problem is convex together with Slater constraint
qualification the KKT-conditions are also sufficient.
\begin{definition}
    Let $x_0 \in X$, where $X$ is the feasible set from equation
    \ref{eq: gen_constraned}. Then $x_0$ fulfills Abadie constraint
    qualification (CQ) if
    \begin{align}
        T_X(x_0) = T_{\text{lin}}(x_0).
    \end{align}
\end{definition}
A particularly nice thing about Abadie CQ is that it does not assume anything
about the objective function $f$ and a very nice result can be proven.
\begin{theorem}
    Let $x^{*}$ be local minimum of \ref{eq: gen_constraned}, fulfilling
    Abadie CQ. Then there exist Lagrange multipliers $\lambda^{*} \in
    \mathbb{R}^{m} $ and $\mu \in \mathbb{R}^{p}$, such that $(x^{*},
    \lambda^{*},\mu^{*})$ is KKT point of the constrained optimization
    problem.
\end{theorem}
\begin{proof}
   The strategy is to utilize the cone duality and already known results.
   From Proposition \ref{prop: descent} we know that $\nabla f(x^{*}) \in
   (T_X(x^{*}))^{*} $, and since $x^{*}$ fulfills Abadie CQ it holds that
   $(T_X(x^{*}))^{*}  = (T_{\text{lin}}(x^{*}))^{*}$. Thus $-\nabla f(x^{*})
   \in -(T_{\text{lin}}(x^{*}))^{*}$. From Theorem \ref{thm: duality} we have
   that $-\nabla f(x^{*}) \in N_{\text{lin} }(x^{*})$ and hence we can write
   \begin{align}
       -\nabla f(x^{*}) = \sum_{i \in \mathcal{A}(x^{*}) } \lambda_i^{*}  \nabla g_i(x^{*}) +
    \sum_{j=1}^{p} \mu_j^{*}h_j(x^{*}),
   \end{align}
   where we know there are $\lambda^{*}_i \ge 0$ for all $i \in
   \mathcal{A}(x^{*})$ and $\lambda^{*}_i = 0$ for all $i \in
   \mathcal{I}(x^{*})$ and $\mu^{*}_j \in \mathbb{R}$ for all $j=1,\ldots,p$.
   Then we expand the sum and rearrange to get
   \begin{align}
       &\nabla f(x^{*}) + \sum_{i=1}^{m} \lambda^{*}_m\nabla g(x^{*})
    + \sum_{j=1}^{p} \mu^{*}_j \nabla h_j(x^{*}) = 0
   \end{align}
   which is the first condition in the KKT conditions
   \begin{align}
       \nabla_x L(x^{*}, \lambda^{*}, \mu^{*}) = 0.
   \end{align}
   For the other conditions, we have for $i$ is active $g_i(x^{*}) = 0$ and
   for $i$ inactive $\lambda_i^{*}=0$. Thus $\lambda^{*}_i \ge 0$,
   $g_i(x^{*})\le 0$,  $\lambda_i g(x^{*}) \le 0$. Also by construction of
   the problem $h_j(x^{*})=0$ for all $j=1,\ldots,p$.
\end{proof}
\section{Convex Case}
Abadie CQ is an excellent result, giving us a hint of finding $x^{*}$.
However, it assumes something about $x^{*}$ before it is even found. In the
case of convex optimization problems we can derive a relaxed global condition
on the optimization problem which implies Abadie CQ. Firstly a brief
introduction on convex functions
\begin{definition}
    Let $U \subseteq \mathbb{R}^{n}$, the function $f: U \to \mathbb{R}$
    is called convex if for all $x, y \in U$ and $\lambda \in [0, 1]$ it
    holds that
    \begin{align}
        f(\lambda x + (1-\lambda)y) \le \lambda f(x) + (1-\lambda)f(y).
    \end{align}
\end{definition}
An important result we will use throughout is the following.
\begin{proposition}
    \label{prop: convexity}
    Let $U \subseteq \mathbb{R}^{n}$ be a nonempty, open and convex set, $f:
    U \to \mathbb{R}$ a differentiable function on $U$. Then the following
    statements are equivalent
    \begin{enumerate}[nosep]
        \item $f$ is convex on $U$;
        \item $\nabla f(x)^{T}(y-x) \le f(y)- f(x)$ for all $x, y \in U$;
        \item $(\nabla f(x) - \nabla f(y))^{T}(y-x)\ge 0$ for all $x, y \in
            U$;
        \item if $f \in C^{2}(U)$, then $\Delta f(x)$ is positive
            semidefinite for all $x \in U$.
    \end{enumerate}
\end{proposition}
\begin{proof}
    We show (1) $\Leftrightarrow$ (2) $\Leftrightarrow$ (3) and (2)
    $\Leftrightarrow$ (4). Starting off with (1) $\Rightarrow$ (2) by definition
    we have that $\forall x, y \in U$ and $\lambda \in [0, 1]$ it holds that
    \begin{align}
        f(\lambda x + (1-\lambda)y) \le \lambda f(x) + (1-\lambda)f(y).
    \end{align}
    Rearranging, dividing by $\lambda$ and expanding the fraction by $x-y$ we
    get
    \begin{align}
        \frac{f(y+\lambda(x-y)) - f(y)}{\lambda (x-y)}(x-y) \le f(x) - f(y).
    \end{align}
    Letting $\lambda \downarrow 0$ we get
    \begin{align}
        \nabla f(y)^{T} (x-y) \le f(x) - f(y).
    \end{align}
For (2) $\Rightarrow$ (1) we know that $U$ is convex hence fix $x, y \in U$
then for all $\lambda \in [0, 1]$ there is a $z$ in the line segment of $x$
and $y$ such that $z = \lambda x + (1-\lambda)y$. We can construct two
inequalities
\begin{align}
    &\nabla f(z)^{T} (x-z) \le f(x) - f(z) \\
    &\nabla f(z)^{T} (y-z) \le f(y) - f(z).
\end{align}
By multiplying the first one with $\lambda$ and second one with $(1-\lambda)$
and adding them together we get
\begin{align}
    \nabla f(z)^{T} (\lambda x + (1-\lambda)y - z) \le \lambda f(x) +
    (1-\lambda)f(y) - f(z).
\end{align}
Inserting for $z$ we get $\lambda x+ (1-\lambda)y - z =0$ and thus
\begin{align}
    f(z) = f(\lambda x + (1-\lambda)y) \le \lambda f(x) + (1-\lambda)f(y).
\end{align}
For (2) $\Rightarrow$ (3) we consider two inequalities for all $x, y \in U$
\begin{align}
    \nabla f(x)^{T} (y-x) \le f(y) - f(x) \\
    \nabla f(y)^{T} (x-y) \le f(x) - f(y).
\end{align}
Adding them together gives
\begin{align}
    (\nabla f(x) - \nabla f(y))^{T} (x-y) \le 0 .
\end{align}
For (3) $\Rightarrow$ (2) we consider a function $g(t) = f(t x + (1-t)y)$
defined on the line segment of $x, y \in U$ for $t \in [0, 1]$,
well-defined since $U$ is convex. The derivative of $g$ is
\begin{align}
    g'(t) = \nabla f(t x + (1-t) y)^{T} (y-x)
\end{align}
Note that by (3) we have that
\begin{align}
    \nabla f(x)^{T} (y-x) \le \nabla f(y)^{T} (y-x),
\end{align}
for all $x, y \in U$. Now we use the Mean Value Theorem on $g$,
\begin{align}
    f(y) - f(x) &= \int_0^1 g'(t) dt\\
                &= \int_0^1 \nabla f(tx +(1-t)y)^{T}(y-x)dt\\
                &\ge \int_0^1 \nabla f(x)^{T} (y-x) dt\\
                &= \nabla f(x)^{T} (y-x).
\end{align}
Lastly for (2) $\Leftrightarrow$ (4) we use Taylor's Theorem, i.e. for all $x, y
\in X$ we have
\begin{align}
    f(y) = f(x) + \nabla f(x)^{T} (y-x) + \frac{1}{2}(y-x)^{T} \nabla^2
    f(\xi)(y-x),
\end{align}
for some $\xi$ in the line segment of $x$ and $y$. For $\Leftarrow$ we have
that the Hessian is positive semidefinite $\frac{1}{2}(y-x)^{T} \nabla^2 f(\xi)(y-x) \ge 0$
thus rearranging gives
\begin{align}
    \nabla f(x)^{T} (y-x) \le f(y) - f(x)
\end{align}
On the other hand we have that $f$ is convex, i.e. $0 \le f(y)-f(x)- \nabla
f(x)^{T} (y-x) $ and thus
\begin{align}
    (y-x)^{T} \nabla^2 f(\xi)(y-x) \ge 0.
\end{align}
Since $x, y$ was arbitrary and $\xi$ in the line segment of $x$ and $y$ we
can conclude that the Hessian is positive semidefinite which concludes the
whole proof.
\end{proof}
Throughout this section consider the following convex problem
\begin{align}
    \min\;&f(x) \label{eq: convex_constrained}\\
    \text{s.t.:}\;&g_i(x) \le 0,\; i=1,\ldots,m \nonumber\\
        &Ax = b,\; \nonumber\\
        &x \in \mathbb{R}^{n}\nonumber,
\end{align}
where $f, g_i$ are convex and continuously differentiable functions for all
$i=1,\ldots,p$ and the function $h$ are affine linear, i.e. $h(x) = Ax-b$,
where $A \in \mathbb{R}^{p\times n}, b\in \mathbb{R}^{p}$. The feasible set
$X$ of problem \ref{eq: convex_constrained} is a convex set in this case. An
important condition for convex problems is the so called Slater constraint
qualification (CQ).
\begin{definition}
    The optimization problem \ref{eq: convex_constrained} fulfills the Slater
    CQ if there exists an $x' \in \mathbb{R}^{n}$ such that
    \begin{align}
        &g_i(x') < 0,\; i=1,\ldots,m \\
        &Ax' = b.
    \end{align}
\end{definition}
Given Slater CQ we can prove the following

\begin{theorem}
    Let $ x^{*} $ be a local minimum of \ref{eq: convex_constrained} and the
    problem fulfills Slater CQ. Then there exist Lagrange multipliers
    $\lambda^{*} \in \mathbb{R}^{m}$ and $\mu^{*} \in \mathbb{R}^{p}$  such
    $(x^{*},\lambda^{*},\mu^{*})$ is a KKT point of
    \ref{eq: convex_constrained}.
\end{theorem}
\begin{proof}
We show that $x^{*}$ fulfills Abadie CQ, i.e. $T_{\text{lin}}(x^{*})
\subseteq T_X(x^{*})$. Let $d \in T_{\text{lin}}(x^{*})$ and choose an $x'
\in X$ satisfying Slater CQ. Set $d' := x' - x^{*}$, by Proposition
\ref{prop: convexity} (2) we have for all $i \in \mathcal{A}(x^{*})$ the
following
\begin{align}
    \nabla g_i(x^{*})^{T} d'
    &= \nabla g_i(x^{*})^{T} (x^{*} + d' - x^{*} )\\
    &\le g_i(x^{*} +d') -g_i(x^{*}) \\
    &< 0,
\end{align}
where $g(x') < 0$ by Slater CQ and $g_i(x^{*}) = 0$ for all active indices
which we considered. Also
\begin{align}
    Ad' = Ax' - Ax^{*} = b - b = 0.
\end{align}
Let $d(\tau) = d+ \tau d'$ for all $\tau >0$, note that $d = \lim_{\tau \to
0}d(\tau)$. Then we have
\begin{align}
    \nabla g_i(x^{*})^{T} d(\tau)
    &= \nabla g_i (x^{*})^{T} d + \tau \nabla g_i(x^{*})^{T} d' < 0,
\end{align}
for all $i \in \mathcal{A}(x^{*})$ and
\begin{align}
    Ad(\tau) = Ad + \tau Ad' = 0.
\end{align}
We show that $d(\tau) \in T_X(x^{*})$ for all $\tau >0 $, which concludes
the proof by the closure of $T_X(x^{*})$. The cone $T_X(x^{*})$ in the case
$X$ is convex can be written as $T_X(x^{*}) = \text{cl}(\text{cone}(X-x^{*}))$. Let
\begin{align}
    x^{k}  = x^{*} + \frac{1}{k}d(\tau),\\
    t_k = 1/k,
\end{align}
for $k \in \mathbb{N}$. This gives
\begin{align}
    \frac{x^{k} - x^{*}  }{t_k} = d(\tau)
\end{align}
To finish the proof, we need to show that $x^{k}  \in X$.
We show that $g_i(x^{k}) \le 0$ for all $i=1,\ldots,m$ and $Ax^{k} -b =
0$ for all $k$ large enough. Indeed for all $i \in \mathcal{A}(x^{*})$ by the
Mean Value Theorem, we can find for each $k\ge 1$ a $\xi^{k}$ in the line
segment of $x^{*}$ and $x^{k}$ such that
\begin{align}
    g_i(x^{k})
    &= g_i(x^{k}) - g_i(x^{*}) \\
    &= \nabla g_i(\xi^{k})^{T} (x^{k} -x^{*} ) \\
    &= \frac{1}{k} \nabla g_i (\xi^{k})^{T} d(\tau).
\end{align}
Thus there is a $k^i_0\ge1$ such that
\begin{align}
    kg_i(x^{k}) < 0 \quad \forall k \ge k_0^{i},
\end{align}
and therefore
\begin{align}
    g_i(x^{k}) < 0 \quad \forall k \ge k_0^{i},
\end{align}
For $i \in \mathcal{I}(x^{*})$, by definition $g_i(x^{*}) < 0$ and $x^{k} \to
x^{*}$ by letting $k$ go to infinity, i.e.
\begin{align}
    \lim_{k \to \infty} g_i(x^{k}) = g_i (x^{*}) < 0.
\end{align}
So there is a $k_1^{i} \ge 1$ such that
\begin{align}
    g_i(x^{k}) < 0 \quad \forall k \ge k_1^{i}.
\end{align}
Taking $k_0 = \max_{i=1,\ldots,m}\{k_0^{i} , k_1^{i}\}$
shows the first part. Lastly
\begin{align}
    Ax^{k} = A x^{*}  + \frac{1}{k}Ad(\tau) = b,
\end{align}
since $x^{*} \in X$ and $Ad(\tau) = 0$. This shows $d(\tau) \in T_X(x^{*})$
and by continuity $d \in T_X(x^{*})$.
\end{proof}
Thus Salter CQ implies Abadie CQ and KKT-conditions in this case
are also sufficient for optimality.
\begin{theorem}
    Let $(x^{*}, \lambda^{*},\mu^{*})$ be a KKT point of
    \ref{eq: convex_constrained}. Then $x^{*}$ is a local minimum of
    \ref{eq: convex_constrained}.
\end{theorem}
\begin{proof}
    Since $x^{*} \in X$ and $(x^{*}, \lambda^{*}, \mu^{*})$ a KKT point of
    \ref{eq: convex_constrained} and by Proposition \ref{prop: convexity} (2)
    we have
    \begin{align}
        f(x)
        &\ge f(x^{*}) + \nabla f(x^{*})^{T}(x-x^{*})\\
        &= f(x^{*}) + \left(
        -\sum_{i=1}^{m} \lambda_i^{*}\nabla g_i(x^{*})^{T}
        -\sum_{j=1}^{p} \mu_j^{*} a_j^{T}
    \right)(x-x^{*})\\
        &= f(x^{*}) -
        -\sum_{i\in\mathcal{A}(x^{*})}^{m} \lambda_i^{*}\nabla g_i(x^{*})^{T}
    (x-x^{*}),
    \end{align}
    where $a_j$ is the $j$-th row of $A$. Using the convexity of $g$ and the
    contraint of $g(x) \le 0$ we get
    \begin{align}
        f(x)
        &\ge f(x^{*})
        -\sum_{i\in\mathcal{A}(x^{*})}^{m} \lambda_i^{*}
        \left( g_i(x) -g_i(x^{*})  \right)\\
        &= f(x^{*})
        -\sum_{i\in\mathcal{A}(x^{*})}^{m} \lambda_i^{*}
        g_i(x)\\
        &\ge f(x^{*}).
    \end{align}
\end{proof}


